\section{Deployment: The Ranger-in-the-Loop Flywheel}
\label{sec:deployment}

The detector ships inside an open-source forest-monitoring platform that we
release with this paper (name withheld for anonymous review).
A field sensor uploads a clip; the API scores it with the fitted detector and,
when $s(x)\ge\tau$, creates an anomaly alert carrying the anomaly score, the
predicted kind, and the full likelihood vector including \emph{unknown}. The
dashboard renders the ranked breakdown (e.g.\ ``likely chainsaw 71\% $\cdot$
vehicle 12\% $\cdot$ unknown 17\%'') so a ranger sees not just \emph{that}
something is unusual but the current best guess and its uncertainty.

\paragraph{Closing the loop without retraining.} Ranger confirmations in the
clip-review UI are exported as a positives manifest and folded back into the
prototype set. Because attribution is closed-form over a frozen embedding, this
is a seconds-long refit, not a training run: the very labels the deployment
generates are the labels that light up its classes. This is the operational
realization of the $k$-curve in Section~\ref{sec:experiments}: each
confirmation moves the system rightward along it.

\paragraph{Why background-first fits social-impact constraints.} Under-resourced
conservation teams cannot commission labeled threat corpora for every reserve,
but every site can record its own background on day one. Anchoring the detector
to background makes the expensive resource (labeled positives) optional and
incremental, and makes the tunable knob a single interpretable
quantity, the background false-positive budget that bounds wasted ranger
patrols. The honest \emph{unknown} option means a novel threat (a new machine, a
new vehicle) is surfaced rather than silently misfiled, which is the failure
mode a fixed taxonomy cannot avoid.

\paragraph{Reproducibility.} The detector math is a small, dependency-light
module unit-tested without the audio stack; the fit and evaluation are single
commands over public manifests; and fitted artifacts, the evaluation protocol,
and the platform are released open-source. We complete the AAAI reproducibility
checklist and document dataset licenses for redistribution.
